% ===================================================================
\chapter{Conclusions and Future Work}
\label{ch:conclusion}

\section{Summary}

This thesis compared data compression methods for SoC-based radar on real,
public data from a radar with many transmitters. It selected the methods with
the best trade-off between compression and loss. It then built their encoders
and decoders on an FPGA and measured them on a board. The four objectives in
Section~\ref{sec:objectives} were met. Eight methods were compared in MATLAB on
50 ColoRadar frames from a cascade radar with twelve transmitters. DRHE and a
second-order LPC were written in C++ for Vitis HLS, together with their
decompressors and a range-FFT front end. These were built into a test system on
an AMD KCU116 board, and checked bit for bit at every step from MATLAB to the
board.

\section{Findings}

\begin{enumerate}
\item Lossless predictive coding gives the best trade-off. DRHE (\num{3.75})
and LPC (\num{3.70}) are the only methods that compress by more than three
without adding a missed detection. The best lossy quantiser reaches
\num{3.76}, but it misses \SI{10.4}{\percent} of the reference detections.
Storing the data as floating point only compresses it by \num{1.04}. All of the
detection loss in the lossless chain comes from the 16-bit quantisation step.
This loss cannot be removed without losing almost all of the compression.

\item On a TDM-MIMO radar, prediction must use the last ramp from the same
transmitter, which is $\nTx$ ramps back. Predicting from the adjacent ramp
compresses worse than not predicting at all. Predicting from the same
transmitter gains \SIrange{6}{8}{\percent}.

\item Most of the ratio comes before prediction. It comes from the unused top
bits of each fixed-point code, and from the many values close to zero. The
dictionary also matters more than the predictor. Dictionaries re-derived from
five frames raise DRHE to \num{4.22} and LPC to \num{4.40} on the board, and
they work just as well on the 45 frames they did not see. At \num{4.22}, the
\SI{0.835}{Gbit/s} from one cascade becomes about \SI{0.20}{Gbit/s}, so four
cascades fit on one Gigabit Ethernet link. At the worst frame (\num{3.5}),
three still fit.

\item The hardware is exact and fast. The output of both compressors was the
same as the simulation, byte for byte, on all 50 frames. The board's decoders
also rebuilt every frame exactly. DRHE compresses a frame in
\SI{0.61}{\milli\second} (\SI{5.1}{Gbit/s}) with \SI{13}{\micro\second} of
latency. LPC takes \SI{3.25}{\milli\second} with \SI{3.0}{\milli\second} of
latency, and both decoders take about \SI{2}{\milli\second}. The encoder and
decoder together use about \SI{23}{\percent} of the chip's logic.

\item DRHE is the better choice for a streaming sensor. It needs no frame
buffer and its cost grows with the data rate. LPC's buffer and latency grow
with the frame size. Kiem's DRHE also works on a cascade radar with only the
prediction lag changed.
\end{enumerate}

\section{Future work}
\label{sec:futurework}

\begin{itemize}
\item \textbf{Fixed-point DRHE.} Kiem's fixed-point types would remove most of
the logic, DSPs and pipeline stages for a \SI{0.07}{\percent} lower ratio, and
the closed loop would keep the design exactly lossless.
\item \textbf{A streaming LPC.} An LPC fitted to previous frames, with no frame
buffer or extra latency.
\item \textbf{Adaptive dictionaries.} Dictionaries trained on more scenes and
updated for each frame.
\item \textbf{All sixteen channels and measured power.} Wider cores for all sixteen
channels, and measured board power.
\item \textbf{The quantisation trade-off.} How many bits of the noise floor are
worth keeping for the detections they save.
\item \textbf{Other MIMO schemes and methods.} Doppler-division and
code-division MIMO, and learned methods such as AdaRadar \cite{park2026}, tested
on the same data.
\end{itemize}

\section*{Use of AI tools}
\addcontentsline{toc}{section}{Use of AI tools}

Generative AI tools were used to help produce the MATLAB figures, the tables
and the formatting of this report. All of the code written for this
thesis was also developed with the support of AI tools.
